(1) \(\nabla^2\Phi=(\Phi_0/r_0^2)(1+1-2)=0\)。\(\mathbf E=-\nabla\Phi=(-\Phi_0 x/r_0^2,-\Phi_0 y/r_0^2,2\Phi_0 z/r_0^2)\)，原点无电荷，其线性场具有零迹梯度。

(2) \(m\ddot x=-q\Phi_0x/r_0^2,\;m\ddot y=-q\Phi_0 y/r_0^2,\;m\ddot z=2q\Phi_0z/r_0^2\)。若 \(q\Phi_0>0\)，横向为稳定谐振，纵向为指数失稳；反号时横向同样失稳，静态电场总无法把所有方向同时稳定。

(3) 记 \(\omega=\sqrt{q\Phi_0/(mr_0^2)}\)，则 \(x=A\cos\omega t+B\sin\omega t\)、\(y=C\cos\omega t+D\sin\omega t\)，\(z=E e^{\sqrt2\omega t}+F e^{-\sqrt2\omega t}\)。真实微小偏离包含增长分支，不能用沿特定稳定流形的精细初态宣称三维俘获。

(4) \(V=q\Phi\) 的Hessian为 \((q\Phi_0/r_0^2)\operatorname{diag}(1,1,-2)\)，迹为零，无法正定。能量 \(H=m(\dot x^2+\dot y^2+\dot z^2)/2+q\Phi\) 守恒但并非下有界。射频Paul阱引入显式时间依赖，可形成动态赝势，已不属于此静态Earnshaw问题。
